\SourceLeaf{249}%
\noindent in abstracto betrachtet nicht widersprechend ist, und füglich höhern Wesen
zukommen könnte. Um aber, aus dieser Höhe, wieder auf die Erde
herunterzukommen, so ist es schade, dass die Gleichheit der Volumina
körperlicher bloss symmetrischer, aber nicht congruenter Gebilde, sich nur
durch die Exhaustionsmethode, und nicht eben so elementarisch demonstriren
lässt, wie meines Wissens zuerst Sie bei der Area des sphärischen Dreiecks
gezeigt haben. Ihr Beweis in etwas anderer Gestalt findet sich auch bei
\textsc{Lobatschewsky}, ohne dass erhellt, ob er Ihren Aufsatz gekannt habe.
Jedes sphärische Dreieck lässt sich nemlich in drei gleichschenklige
Dreiecke zerlegen, und diese Zerlegung auf zwei symmetrische Dreiecke
angewandt, gibt bei beiden gleiche und nur in ungleicher Ordnung liegende
Dreiecke. (Es versteht sich, dass um vollständig zu sein, man drei Fälle
unterscheiden müsste, wovon der zweite und dritte auch so ausgesprochen
werden können, dass je eins der Partialdreiecke negativ oder 0 sein kann.)
\mbox{.\ .\ .\ .\ .\ .\ .\ .}

\par\bigskip

\Needspace{7\baselineskip}
\begin{center}
[\textsc{Gerling} \emph{an} \textsc{Gauss}.\quad \emph{Marburg, 15. April 1844.}]
\end{center}

\{Obwohl ich Ihren Brief vom 8. d. M., den ich vorläufig herzlich verdanke,
eigentlich noch gar nicht beantworten kann, weil ich bis jetzt noch nicht
dazu habe kommen können, ihn so zu studiren und in mich aufzunehmen, als
sein reichhaltiger Inhalt und mein Interesse für den Gegenstand erfordert;
so kann ich mir doch das Vergnügen nicht versagen, Ihnen in Beziehung auf
eine einzelne Äusserung darin ein Paar Zeilen zu schreiben.

Sie sagen: »Es ist schade, dass die Gleichheit der Volumina körperlicher
bloss symmetrischer, aber nicht congruenter Gebilde sich nur durch die
Exhaustionsmethode, und nicht eben so elementarisch demonstriren lässt«
etc.

Erlauben Sie, dass ich Ihnen hier eine solche Demonstration mittheile mit
der Bitte, mir gelegentlich Ihr Urtheil darüber zu schreiben.

\textls{Lehrsatz}: Zwei beliebige symmetrische Polyeder sind von gleichem
körperlichen Inhalt.

\textls{Beweis}: I. Dieselben lassen sich durch angemessen durchgelegte Ebenen
